\documentclass[11pt,a4paper]{article}
\usepackage[margin=22mm]{geometry}
\usepackage[T1]{fontenc}
\usepackage{lmodern,microtype,graphicx}
\usepackage{xurl,fancyvrb,upquote}
\usepackage[hidelinks]{hyperref}
\DefineVerbatimEnvironment{commands}{Verbatim}{fontsize=\footnotesize,samepage=true}
\setlength{\parindent}{0pt}
\setlength{\parskip}{4pt}
\title{FFA26 workshop\\\large OPALX installation with Docker}
\author{}
\date{16 November 2026}

\begin{document}
\noindent
\begin{minipage}[c]{0.68\textwidth}
\centering
{\LARGE FFA26 workshop}\\[4pt]
{\large OPALX installation with Docker}\\[6pt]
16 November 2026
\end{minipage}\hfill
\begin{minipage}[c]{0.28\textwidth}
\includegraphics[width=\linewidth]{opalx-logo.png}
\end{minipage}
\par\vspace{8pt}

The workshop uses \textbf{OPALX in a Docker container} for simulation and
\textbf{Python/Jupyter on your computer} for the School notebook and plots.
Install both before the workshop and complete the checks below. You do not
need to compile OPALX or install MPI locally. These instructions follow the
\href{https://opalx-project.github.io/opalx-manual/getting-started/index.html\#run-with-docker}{current OPALX Docker guide}.

\section{Install and start Docker}
\begin{itemize}
\item \textbf{macOS:} install
\href{https://docs.docker.com/desktop/setup/install/mac-install/}{Docker Desktop for Mac},
choosing Apple silicon or Intel to match your machine. Start Docker Desktop.
If a separate Docker client cannot find \texttt{docker-credential-desktop},
add the bundled tools to this terminal's path:
\begin{commands}
export PATH="/Applications/Docker.app/Contents/Resources/bin:$PATH"
\end{commands}
\item \textbf{Windows:} install
\href{https://docs.docker.com/desktop/setup/install/windows-install/}{Docker Desktop with the WSL 2 backend}
and enable its \href{https://docs.docker.com/desktop/features/wsl/}{Ubuntu WSL integration}.
Run all commands below in the \emph{Ubuntu/WSL terminal}, including Python
and Jupyter. Keep the course in your Linux home directory.
\item \textbf{Linux:} follow the
\href{https://docs.docker.com/engine/install/}{installation} and
\href{https://docs.docker.com/engine/install/linux-postinstall/}{post-installation instructions}
to run Docker as your normal user without \texttt{sudo}.
\end{itemize}

The remaining commands use Bash or Zsh: Terminal on macOS/Linux or Ubuntu
on Windows. Check that Docker can contact its running server:
\begin{commands}
docker version
docker run --rm hello-world
\end{commands}
The first command should show both Client and Server information; the second
should print a successful test message.

\section{Download the OPALX CPU image}
\begin{commands}
export OPALX_IMAGE=ghcr.io/opalx-project/opalx:cpu
docker pull "$OPALX_IMAGE"
docker run --rm "$OPALX_IMAGE" --version-full --info 5
docker run --rm "$OPALX_IMAGE" --help
\end{commands}
This public CPU/OpenMP image supports Linux AMD64 and ARM64; Docker selects
the architecture. No GitHub login is needed. To check the selected architecture:
\begin{commands}
docker run --rm --entrypoint uname "$OPALX_IMAGE" -m
\end{commands}
Expect \texttt{aarch64} on Apple silicon and \texttt{x86\_64} on Intel/AMD.
Do not force \texttt{linux/amd64} on an Apple silicon Mac. If necessary, run
\texttt{unset DOCKER\_DEFAULT\_PLATFORM}, then pull again.

The \texttt{cpu} tag follows successful development builds; it is not a fixed
release. Use the workshop image digest supplied by the organisers when available.
To select the exact image just downloaded:
\begin{commands}
export OPALX_IMAGE="$(docker image inspect \
  --format '{{index .RepoDigests 0}}' "$OPALX_IMAGE")"
printf '%s\n' "$OPALX_IMAGE"
\end{commands}

\clearpage
\section{Prepare the course and Python environment}
Obtain the workshop checkout or ZIP supplied by the organisers. It must contain
\path{opalx/FFA26_OPALX.ipynb}, \path{opalx/pyfiles/}, and the three machine
input directories. Keep this structure intact. In the commands below, replace
\path{/path/to/FFA25-school} with your course directory.
\begin{commands}
cd /path/to/FFA25-school
export FFA_SCHOOL="$(pwd -P)"
python3 --version
\end{commands}
Use \textbf{Python 3.10 or newer}. On macOS, install a current Python from
\href{https://www.python.org/downloads/macos/}{python.org} if needed. On
Ubuntu/WSL, the distribution packages \texttt{python3}, \texttt{python3-venv}
and \texttt{python3-pip} provide the starting environment; check the version.
Create a separate environment and install the notebook packages:
\begin{commands}
python3 -m venv "$HOME/.venv-ffa26"
source "$HOME/.venv-ffa26/bin/activate"
python -m pip install --upgrade pip
python -m pip install numpy pandas matplotlib h5py \
  jupyterlab ipykernel
\end{commands}
The Python environment runs locally, not inside the OPALX runtime container.
All commands in the next section use this same activated environment.

\section{Connect the notebook to Docker}
The notebook expects an executable named by \texttt{OPALX\_EXE}. Create a small
launcher in the Python environment; it forwards each OPALX invocation to Docker:
\begin{commands}
cat > "$VIRTUAL_ENV/bin/opalx-docker" <<'SH'
#!/bin/sh
set -eu
: "${FFA_SCHOOL:?Set FFA_SCHOOL to the course root}"
: "${OPALX_IMAGE:?Set OPALX_IMAGE to the workshop image}"
FFA_RUN_DIR="$(pwd -P)"
case "$FFA_RUN_DIR/" in
  "$FFA_SCHOOL/"*) ;;
  *) echo "Run from inside FFA_SCHOOL" >&2; exit 2 ;;
esac
exec docker run --rm \
  --user "$(id -u):$(id -g)" \
  -e OPALX_NP=1 -e OMP_NUM_THREADS=1 \
  -e OMP_PROC_BIND=false \
  --mount "type=bind,src=$FFA_SCHOOL,dst=$FFA_SCHOOL" \
  --workdir "$FFA_RUN_DIR" \
  "$OPALX_IMAGE" "$@"
SH
chmod +x "$VIRTUAL_ENV/bin/opalx-docker"
export OPALX_EXE="$VIRTUAL_ENV/bin/opalx-docker"
unset OPALX_LAUNCHER OPALX_MPI_FLAGS
"$OPALX_EXE" --version-full --info 5
\end{commands}
The launcher uses one MPI rank and one OpenMP thread. It mounts the whole
course at the \emph{same absolute path} inside the container and keeps the
current run directory. This matters because the notebook writes absolute
field-map output paths. Results are therefore available to local Python and
remain on your computer after the container exits. The user/group option
keeps output ownership consistent on Linux.

\clearpage
\section{Check one ring, then open the notebook}
Before running the lessons, check a standalone native closed-orbit calculation:
\begin{commands}
cd "$FFA_SCHOOL/opalx/school-16-cell"
"$OPALX_EXE" school-16-cell.in --info 1
\end{commands}
OPALX should report a verified closed orbit, a stable one-turn eigenanalysis,
and \texttt{Full tunes (uncoupled planar)}. For the unchanged 12 MeV input,
expect approximately $Q_r=3.41066$ and $Q_Y=3.39111$. The generated
\texttt{orbit.json} stays beside this input. If full tunes are missing, check
the image revision with the organisers before continuing.

Start Jupyter from the same terminal, retaining the exported variables:
\begin{commands}
cd "$FFA_SCHOOL/opalx"
python -m jupyterlab FFA26_OPALX.ipynb
\end{commands}
Open the displayed local URL in your browser and use the Python kernel from
\texttt{.venv-ffa26}. Leave \texttt{EXECUTABLE = None} in the notebook so it uses
\texttt{OPALX\_EXE}. Run from the top: original ring, 55-degree/D-off variant,
then fourfold ring. Each input supplies its validated numerical settings.
The notebook creates a fresh directory under \path{opalx/runs/}; its new COF
JSON supplies the closed state and full tunes used by the tracking lesson.
It does not load the standalone JSON snapshot.

\subsection*{Returning later}
Start Docker, activate the environment, and repeat these exports before
starting Jupyter. Replace the image reference with your recorded workshop
reference if the organisers specify a digest:
\begin{commands}
source "$HOME/.venv-ffa26/bin/activate"
cd /path/to/FFA25-school
export FFA_SCHOOL="$(pwd -P)"
export OPALX_IMAGE=ghcr.io/opalx-project/opalx:cpu
export OPALX_EXE="$VIRTUAL_ENV/bin/opalx-docker"
unset OPALX_LAUNCHER OPALX_MPI_FLAGS
cd opalx
python -m jupyterlab FFA26_OPALX.ipynb
\end{commands}

\section{If a check fails}
\begin{itemize}
\item \textbf{Docker server unavailable:} start Docker Desktop or the Linux
Docker service, then repeat \texttt{docker version}.
\item \textbf{Mount denied or output path missing:} allow Docker Desktop to
share the course directory. Check that \texttt{FFA\_SCHOOL} is the course root
and that you are using the launcher above, rather than mounting only a single
run folder at \texttt{/data}.
\item \textbf{Notebook cannot find OPALX:} check
\texttt{echo "\$OPALX\_EXE"} and rerun the launcher's version command.
Start Jupyter from the terminal containing the exports. After changing the
launcher environment, stop Jupyter and start it again from that terminal.
\item \textbf{Python imports fail:} activate the workshop environment, start
its Jupyter, and restart the notebook kernel.
\end{itemize}
For unresolved problems, retain the Docker image reference, the version/revision
output and the failing command or notebook \texttt{run.log}. Consult the
\href{https://opalx-project.github.io/opalx-manual/}{OPALX manual} or contact the
workshop organisers before the workshop.

\end{document}
